
\section{Plano horizontal}
\label{sec:plano_horizontal}

Após calcular a distância frontal $d$ e a altura $h$ no plano lateral, é necessário considerar a rotação da junta J1 para determinar a posição tridimensional do centro de preensão.

A junta J1 roda o braço em torno do eixo vertical $z$, alterando a direção do seu movimento no plano horizontal $(x,y)$, como ilustrado na Figura~\ref{fig:topo}.

\begin{figure}[htbp]
    \centering
    \resizebox{0.6\textwidth}{!}{%Vista de cima do robô
\begin{tikzpicture}[scale=0.3,
                    % Estilos
                    link/.style={line width=2pt, line cap=round},
                    cota/.style={gray!70!black, thin}]

  % Coordenadas das juntas
    \coordinate (J1) at (0,0);
    \coordinate (J6) at (-4.7,21);     % centro de preensão (x, y)
    \coordinate (O5) at (5.6,3.2);

  % Base
    \draw[->] (-10,0) -- (12,0) node[below left] {$x$};
    \draw[->] (0,-3) -- (0,25) node[left] {$y$};
    \fill[gray!25] (0,0) circle (2.5);  \draw (0,0) circle (2.5);
    \node[below left] at (-2.2,-0.2) {J1};
  
  % Desvio lateral L5 e plano do braço
    \draw[link] (J1) -- node[pos=0.55, below right=-2pt] {$L_5$} (O5);
    \draw[link] (O5) -- node[pos=0.4, right=2pt] {$d$} (J6);
    \draw[cota, dashed] (O5) -- ($(O5)!1.1!(J6)$) node[above, black] {plano do braço};
    \draw[cota, dashed] (O5) -- (30:14);
    \fill (J6) circle (7pt); \node[left=4pt] at (J6) {$(x, y)$};
  
  % Distância
    \draw[cota, dashed] (J1) -- node[pos=0.8, left=1pt, black] {$r$} (J6);
  
  % Ângulos
    \draw[->, thick] (5,0) arc[start angle=0, end angle=30, radius=5] node[pos=0, below right=2pt and -3pt] {$q_1 > 0$};
    \draw[->, thin] (8.5,0) arc[start angle=0, end angle=102.6, radius=8.5] node[pos=0.27, right=2pt] {$\operatorname{atan2}(y, x)$};
    \draw[->, thin] (30:12) arc[start angle=30, end angle=102.6, radius=12] node[pos=0.3, above right=0pt] {$\operatorname{atan2}(d, L_5)$};
    \draw[thin] ($(O5)+(120:1.3)$) -- ($(O5)+(120:1.3)+(210:1.3)$) -- ($(O5)+(210:1.3)$);

  % Convenção de sinal (visto de cima)
  %  \draw[->, thick, red!70!black] (9,23) ++(-30:1.6) arc[start angle=-30, end angle=250, radius=1.6];
  %  \node[red!70!black] at (9,23) {$+$};
  %  \node[right=2pt, align=left, font=\small] at (10.8,23) {ângulo positivo:\\sentido anti-horário};

\end{tikzpicture}}
    \caption{Vista superior do manipulador no plano horizontal.}
    \label{fig:topo}
\end{figure}

\subsection{Desvio lateral}

O manipulador robótico apresenta um desvio lateral $L_5$ de aproximadamente $6{,}5$ cm, devido à disposição dos seus elos.

Isto significa que o plano de movimento do braço não passa exatamente pelo eixo de rotação da base, encontrando-se deslocado lateralmente, como representado na Figura~\ref{fig:topo}.

Na formulação desenvolvida, este desvio é representado por $l$, que inclui também o efeito da rotação da junta J5, conforme definido na Secção~\ref{sec:plano_lateral}.

Assim, a posição horizontal do centro de preensão depende de duas distâncias perpendiculares:
\begin{itemize}
    \item $d$: distância frontal, calculada no plano lateral;
    \item $l$: desvio lateral relativamente ao eixo da base.
\end{itemize}

\subsection{Rotação da base}

Na configuração de referência, quando $q_1=0^\circ$, a distância frontal $d$ encontra-se orientada segundo o eixo positivo $y$, enquanto o desvio lateral $l$ está orientado segundo o eixo positivo $x$.

Por isso, nesta configuração, as coordenadas do centro de preensão são:

\begin{equation}
    x=l,\qquad y=d,\qquad z=h
\end{equation}

Quando J1 roda de um ângulo $q_1$, as direções de $d$ e $l$ também rodam em torno do eixo vertical $z$.

Para determinar as novas coordenadas, projetam-se estas duas distâncias nos eixos $x$ e $y$.

A componente lateral $l$ contribui com $l\cos q_1$ para $x$ e $l\sin q_1$ para $y$.

Por sua vez, a componente frontal $d$, perpendicular a $l$, contribui com $-d\sin q_1$ para $x$ e $d\cos q_1$ para $y$.

Somando estas contribuições, obtêm-se as coordenadas tridimensionais do centro de preensão:

\begin{equation}
    \boxed{
    \begin{aligned}
        x &= l\cos q_1-d\sin q_1\\
        y &= l\sin q_1+d\cos q_1\\
        z &= h
    \end{aligned}
    }
    \label{eq:xyz}
\end{equation}

A coordenada $z$ mantém-se igual a $h$, porque a rotação de J1 ocorre em torno do eixo vertical e não altera a altura da garra.

Na ausência de desvio lateral ($l=0$), as expressões simplificam-se para:

\begin{equation}
    x=-d\sin q_1,\qquad
    y=d\cos q_1,\qquad
    z=h
\end{equation}

Neste caso, o plano de movimento do braço passa pelo eixo de rotação da base.

\subsection{Distância ao eixo da base}

Conhecidas as coordenadas $x$ e $y$, é possível determinar a distância horizontal $r$ entre o centro de preensão e o eixo da base.

Como as distâncias $d$ e $l$ são perpendiculares, formam um triângulo retângulo, pelo que se aplica o teorema de Pitágoras:

\begin{equation}
    r=\sqrt{x^2+y^2}=\sqrt{l^2+d^2}
    \label{eq:r}
\end{equation}

Esta relação permite determinar a distância frontal a partir de uma posição pretendida, sendo utilizada posteriormente na cinemática inversa.

Além da distância, é necessário conhecer a direção do alvo no plano horizontal.

O ângulo formado pela posição $(x,y)$ relativamente ao eixo $x$ é dado por $\operatorname{atan2}(y,x)$.

Este ângulo corresponde à soma da rotação da base $q_1$ com o ângulo formado pelas componentes $l$ e $d$:

\begin{equation}
    \operatorname{atan2}(y,x)
    =
    q_1+\operatorname{atan2}(d,l)
    \label{eq:angulo_horizontal}
\end{equation}

Esta expressão será utilizada na cinemática inversa para determinar o ângulo $q_1$ necessário para orientar o manipulador na direção de um alvo.
